\subsection{F(h): local stratifications, flags, and controlled links}

The topology used in this subsection is essential. A stratum has its
subspace topology in a global chart; a global leaf of a local
stratification has the topology obtained by gluing its local plaques.
A neighbourhood germ remains a pro-object of \emph{full} sheaf topoi.
Neither taking connected components nor taking covering categories is
part of that definition.

\paragraph{F.31 (Proved). With a counterexample to the sheaf assertion:
the presheaf is separated, and its sheafification is the atlas of local
stratifications.}
Use conditions a)--d) of MPL~147, fol.~86, scans~181--182: a locally
finite partition into locally closed, connected, locally connected strata;
the frontier condition; and, for every stratum $S$ and $x\in\overline S$,
a fundamental system of neighbourhoods $N$ of $x$ for which $N\cap S$
is nonempty and connected. Restrictions split strata into connected
components. These restrictions define a separated presheaf
$\operatorname{Strat}_{\rm glob}$. It is not a sheaf, even if all the
charts satisfy e). A section of its associated sheaf is exactly an open
atlas of such partitions agreeing, after component splitting, on
overlaps. Its leaves and all its local flag spaces are intrinsic to that
section, rather than to a chosen atlas.

\begin{proof}
We first verify the restriction assertion, which is needed before asking
whether the presheaf is separated. Let $U\subset X$ be open. A component
$C$ of $S\cap U$ is open in $S\cap U$ because $S$ is locally connected,
and is therefore locally closed in $U$. For $x\in U$, only finitely many
original strata meet a sufficiently small neighbourhood. For each such
stratum whose closure contains $x$, choose $N_S\subset U$ with
$N_S\cap S$ connected. This set lies in just one component of $S\cap U$.
Intersect the finitely many $N_S$, and avoid the closures of the remaining
strata. The resulting neighbourhood meets at most one component of each
original stratum, proving local finiteness. If $x\in\overline C^{,U}$,
every sufficiently small connected intersection $N\cap S$, with
$N\subset U$, meets $C$ and is contained in it. This proves condition d)
for $C$.

Let $C\subset S\cap U$ and $E\subset R\cap U$ be components and suppose
$x\in E\cap\overline C^{,U}$. The original frontier condition gives
$R\subset\overline S$. Choose $N\subset U$ about $x$ with $N\cap S$
connected, hence contained in $C$. For every $y\in E\cap N$, every
neighbourhood of $y$ inside $N$ meets $S$, hence $C$. Thus
$E\cap\overline C^{,U}$ is open in $E$; it is also closed and nonempty.
Connectedness of $E$ gives $E\subset\overline C^{,U}$. This is the
frontier condition for the restricted partition. Successive component
splitting agrees with direct restriction. Condition e), when imposed as
a condition on local product germs, is preserved by the same restriction.

Suppose global partitions $P,Q$ have the same restriction to each member
of an open cover. Fix a stratum $S$ of $P$. Its intersections with the
cover split into connected plaques, open in $S$. Each plaque is also a
plaque for $Q$. Two overlapping plaques belong to the same stratum of
$Q$. The graph whose vertices are these plaques and whose edges are
nonempty intersections is connected: a partition into two graph
components would partition $S$ into two nonempty disjoint open subsets.
Consequently $S$ lies in one $Q$-stratum. Interchanging $P,Q$ proves
equality of the partitions. This proves separatedness, answering the
margin of scan~182 affirmatively. It also explains why the deleted word
``separated'' cannot be restored in the preceding sentence.

For failure of gluing, choose distinct $p,q\in S^1$ and put
$U=S^1\setminus\{q\}$, $V=S^1\setminus\{p\}$. Stratify $U$ by $p$ and
the two components of $U\setminus\{p\}$, and stratify $V$ by its single
stratum. Both restrictions to $U\cap V$ are its two connected
components. Both charts satisfy a)--e). A global partition with these
restrictions must have strata $\{p\}$ and $S^1\setminus\{p\}$: the
second follows from its connectedness and the preceding plaque argument.
The germ of $S^1\setminus\{p\}$ at $p$ has two components, violating d).
This is also the boundary phenomenon in the marked-disc picture of
scan~181.

Because the presheaf is separated, its sheafification is its usual
``plus'' construction: a compatible atlas modulo common refinement.
To construct the leaves, generate an equivalence relation on $X$ by
membership in a connected local plaque. A leaf is given the final
topology for all its plaques. Plaques agree on open subplaques, so they
are an atlas of locally closed embeddings into $X$. The leaf inclusion
is a local immersion, and the disjoint union of leaves is the bijective
local immersion $D_0^*\to X$ of (296)--(297). This construction commutes
with open restriction. If $X$ is compact Hausdorff, a finite collection
of charts, each reduced to finitely many plaques, covers $X$; every leaf
contains one of these plaques. There are therefore finitely many leaves,
and they are locally compact. The Hausdorff condition is included here
to justify the local compactness assertion of scan~184.
\end{proof}

\paragraph{F.32 (Proved). The germ flags form the nerve of an internal
topological category in every degree.}
For a chart $V$ with strata $S_i$, order them by
$S_i\leq S_j\iff\overline S_i^{,V}\subset\overline S_j^{,V}$.
The flag spaces of (299)--(303), with the literal topology specified
there, are
\[
 D_r|_V=
 \coprod_{i_0\leq\cdots\leq i_r}\overline S_{i_0}^{,V},
 \qquad
 D_r^*|_V=
 \coprod_{i_0\leq\cdots\leq i_r}S_{i_0}.
 \tag{F.h.1}
\]
These formulas glue canonically for every local stratification of F.31.
Every projection $D_r\to X$ is a local immersion; $D_r^*\subset D_r$
is open. Deleting or repeating entries gives continuous face and
degeneracy maps. For every $r\geq2$ the Segal map is a homeomorphism
\[
 D_r\ \xrightarrow{\ \sim\ }\
 \underbrace{D_1\times_{D_0}\cdots\times_{D_0}D_1}_{r\text{ factors}}.
 \tag{F.h.2}
\]
Thus $D_\bullet$ is exactly the nerve of an internal category in
topological spaces. Its fibre over $x$ is the nerve of the finite poset
of incident stratum germs at $x$. The first-vertex map $D_r\to D_0$ is
a finite covering locally on $D_0$; in particular it has the asserted
property over $D_0^*$. Under the literal model (F.h.1), this property
over all of $D_0$ is not special to normal crossings.

\begin{proof}
Let $W\subset V$ be open and $x\in\overline S_i\cap W$. Condition d)
selects a unique component $C_i$ of $S_i\cap W$ whose closure contains
$x$: a small connected $N\cap S_i$ lies in such a component, and any
other component accumulating at $x$ would meet that same $N$.
The choice is locally constant as $x$ moves in $\overline S_i\cap W$.
Indeed $N\cap S_i\subset C_i$ implies
$N\cap\overline S_i\subset\overline C_i$.
For a chain $S_{i_0}\leq\cdots\leq S_{i_r}$ these uniquely selected
components form a chain: near $x$, the lower component is contained in
the closure of the higher one, and the restricted frontier condition
proved in F.31 extends this inclusion over the entire connected lower
component. This gives a canonical homeomorphism between the restriction
of (F.h.1) and the formula computed after component splitting on $W$.
The identifications compose exactly, because they select the same germs.
They therefore glue on an arbitrary atlas and prove independence of
that atlas. This also gives the precise sense in which the associated
sheaf is described by local flags: an atlas defines compatible germ
flags, and the degree-zero strict pieces recover its local plaques.
One must retain the topology and the maps to $X$; bare incidence posets
do not determine a local stratification.

Deleting the first entry of a chain uses the inclusion
$\overline S_{i_0}\hookrightarrow\overline S_{i_1}$; deleting any other
entry uses the identity on the first closure. Repetition also uses the
identity. All are continuous and commute with chart changes. The
simplicial identities are the identities for deleting and repeating
entries, before any topological quotient is taken.

In the iterated fibre product in (F.h.2), composable arrows have the
same origin $x$ and the same middle \emph{germ}. Their concatenation
is a unique chain. On the component indexed by this chain, the fibre
product is
\[
 \overline S_{i_0}
 \times_{\overline S_{i_1}}\overline S_{i_1}
 \times_{\overline S_{i_2}}\cdots
 \times_{\overline S_{i_{r-1}}}\overline S_{i_{r-1}}
 \cong\overline S_{i_0}
\]
with its subspace topology. This proves a homeomorphism, not just a
bijection, and works for every $r$. Composition is concatenation followed
by forgetting the intermediate entry; associativity and units follow
from the same identities. The strict spaces alone do not form this
category, since deleting the initial stratum need not preserve strictness.

Fix a chart stratum $S_i$. Every stratum above it has closure containing
any selected point of $S_i$, so local finiteness implies that there are
only finitely many such strata. Over the copy of $\overline S_i$ in
$D_0$, the first-vertex map is a disjoint union of a finite number of
identity maps, one for every chain beginning with $i$. This proves the
covering assertion, including at boundary points of that copy. Thus the
aside in scan~189 claiming that extension from $D_0^*$ to all $D_0$ is
special to normal crossings is not a consequence of the definitions
as written. The manifolds and the Boolean incident posets of normal
crossings \emph{are} additional properties. Likewise the phrase about
degeneracies giving isomorphisms at scan~187 means isomorphisms onto the
equal-adjacent-entry summands, not $D_r\cong D_{r+1}$: the two-element
chain already has two vertices and three arrows. Passing to $\pi_0$
before the fibre products would invalidate this proof; F.25 addresses
that separate error.
\end{proof}

\paragraph{F.33 (Refuted). Product germs in e) do not imply a stationary
link or a conical tube.}
There is a compact metrizable space with two strata satisfying a)--e)
whose punctured neighbourhood pro-homology at the closed stratum is not
constant. Hence e) alone cannot justify replacing its elementary
neighbourhood pro-topoi by fixed link models. This does not obstruct
the full pro-topos descent theorem below.

\begin{proof}
Construct an orientable surface $M$ by starting with a disc and attaching
successively compact genus-one cobordisms with two boundary circles.
Choose a compact exhaustion $K_n$ by connected genus-$n$ subsurfaces
with one boundary circle and $K_n\subset\operatorname{int}K_{n+1}$.
Then $M$ has one end and infinite genus. Let
$X=M\cup\{\infty\}$ be its one-point compactification. It is compact
metrizable and is stratified by $M$ and $\{\infty\}$. Both strata are
connected and locally connected; the partition is finite and satisfies
the frontier condition. The sets
$T_n=M\setminus K_n$ are a connected cofinal system of punctured
neighbourhoods of $\infty$, proving d). Along $M$, e) is the product
with a point. Along the point $\infty$, e) is the tautological product
$\{\infty\}\times(X,\infty)$. Its transverse stratum has connected
germ. Thus even the asserted transverse connectedness in scan~182 holds.

Choose the two handle generators $a_j,b_j$ in the $j$th cobordism and
let $c_n$ be its inner boundary class. Cutting off after the $m$th
handle gives a genus-$(m-n)$ surface with two boundary components,
whose first homology has basis
$c_n,a_{n+1},b_{n+1},\ldots,a_m,b_m$. Passing to the union yields
\[
 H_1(T_n;\mathbb Z)
  =\mathbb Z c_n\oplus\bigoplus_{j>n}
       (\mathbb Z a_j\oplus\mathbb Z b_j).
\]
Inclusion $T_m\hookrightarrow T_n$, $m>n$, sends $c_m$ to $c_n$
and preserves the surviving handle generators. Its image is
\[
 \mathbb Z c_n\oplus\bigoplus_{j>m}
       (\mathbb Z a_j\oplus\mathbb Z b_j),
 \tag{F.h.3}
\]
which strictly decreases with every $m$.

For completeness, an inverse system isomorphic in pro-abelian groups
to a constant group has eventually stationary images at each fixed
stage. Write such inverse isomorphisms as
$f:A_n\to B$ and compatible $g_m:B\to A_m$ after choosing a
representative of $f$. The equation $gf=1$ says that for sufficiently
large $k$ the transition $A_k\to A_m$ factors through $g_m$; the
compatibility $g_m=(A_k\to A_m)g_k$ gives the reverse inclusion of
images. Hence the eventual image equals $g_m(B)$. Formula (F.h.3)
contradicts this conclusion. A conical punctured tube
$L\times(0,\epsilon)$ with compatible shrinkings would have constant
pro-homology $H_1(L)$, so no such link model exists here. A finite
truncation of the image calculation is reproduced in the accompanying
program; the argument above proves strict decrease at every stage.
\end{proof}

\paragraph{F.34 (Proved). Elementary pro-topoi, their gluing, and the
precise additional link hypotheses.}
There are two different assertions.
First, full sheaf descent holds for neighbourhood \emph{pro-topoi}
without a conical or asphericity assumption. In an ambient open $Y$
where a locally closed set $Z$ is closed, define
\[
 \mathfrak T_{Z,Y}=(\operatorname{Sh}(N))_{N\supset Z},
 \qquad
 \mathfrak L_{Z,Y;W}=(\operatorname{Sh}(N\cap W))_{N\supset Z},
 \tag{F.h.4}
\]
where $W$ is a specified open part of $Y$; in particular
$W=Y\setminus Z$ gives the punctured tube. Restriction to local
stratification charts and all their overlaps defines the same objects
for local plaques. For an incident flag
$\sigma=(S_0<\cdots<S_r)$, its elementary pro-topoi are the iterated
neighbourhood objects
\[
 \mathfrak E_{\sigma,W}
   =\bigl(\operatorname{Sh}
       (W\cap N_0\cap\cdots\cap N_r)\bigr)_{(N_0,\ldots,N_r)},
 \tag{F.h.5}
\]
where $N_j$ ranges over neighbourhoods of the indicated plaque $S_j$
in the chart, and $W$ is an independently specified locally closed
union of strata, with its \emph{subspace topology}. It need not be
open in the ambient chart. The empty flag gives the constant topos
$\operatorname{Sh}(W)$. One may take $W=S_r$, in which case the last
neighbourhood condition is redundant, or take a stratum closure, or
the specified punctured part of such a closure. Ambient open
restrictions make every closed embedding in this expression legitimate.
In a Cech or flag diagram $W$ is held fixed. Refining neighbourhoods
and deleting a neighbourhood condition give the comparison maps;
a specified inclusion $W\subset W'$ gives a further comparison.
There is no inclusion from one disjoint stratum into another when a
terminal vertex is deleted. Conical retractions, when available, are
separate geometric maps and retain their specified control data.
Common refinement makes the restriction identities strict. These
definitions do not assert that
$\mathfrak E_{\sigma,W}$ is constant.

Second, to model these objects by locally trivial link bundles, it is
sufficient to assume a locally compact paracompact Hausdorff space of
finite depth with compatible conical control data: local manifold
plaques $S$, tubes $T_S$, retractions $\pi_S$, and radius functions
$\rho_S$ with $\rho_S^{-1}(0)=S$, locally identified over $S$ with
\[
 S\times C(L_S),\qquad
 T_S\setminus S\cong L_S^{\rm bun}\times_S(S\times(0,\epsilon_S)).
 \tag{F.h.6}
\]
Here $L_S$ is compact and recursively locally contractible and
conically stratified; $L_S^{\rm bun}\to S$ is a locally trivial bundle,
so monodromy is retained. Changes of charts preserve the strata, cone
radius, and these control data. For incident plaques $S<T$ require
$\pi_S\pi_T=\pi_S$ and $\rho_S\pi_T=\rho_S$ wherever defined;
tubes of incomparable plaques are disjoint after shrinking. The data
may be given on a locally finite atlas with its full overlap nerve,
which allows distinct local branches of one global leaf. This is an
explicit sufficient hypothesis, not a conclusion from e). Whitney
stratifications in the usual locally compact smooth setting provide
such controlled tube germs by the control-data and isotopy results of
Mather (1970, \S7 and Corollaries~10.2--10.6)\cite{Mather1970Local}.

With these data, the full Cech diagram of tubes covering a closed
union of strata, or its subdivision by all flags, has the required
two-colimit in pro-Grothendieck topoi. In particular, for a closed
union $Z\subset Y$,
\[
 \operatorname{Sh}(Y)\ \simeq\
 \operatorname{Sh}(Y\setminus Z)
       \amalg_{\mathfrak L_{Z,Y;Y\setminus Z}}
       \mathfrak T_{Z,Y}
 \quad\text{in pro-Grothendieck topoi}.
 \tag{F.h.7}
\]
At the level of homotopy types these are homotopy colimits of the
\emph{full} link and iterated-link bundles. No $K(\pi,1)$ assumption
is required for this gluing. To replace every elementary homotopy type
by its fundamental groupoid, the exact condition is that every such
type be a $1$-type. A sufficient geometric condition is that the base
strata and all the relevant iterated-link fibres be $1$-types; for
connected components this means $K(\pi,1)$ spaces. The reconstruction
then uses the homotopy colimit of their classifying \emph{spaces}, not
an ordinary two-colimit of classifying topoi or groupoids.

\begin{proof}
At a fixed open-neighbourhood stage, an open cover $N=\bigcup N_a$
gives the usual sheaf descent equivalence, including morphisms:
a sheaf on $N$ is a family of sheaves on $N_a$ with isomorphisms on
$N_a\cap N_b$ satisfying the triple cocycle. For any target topos
$\mathcal A$, apply this objectwise to inverse-image functors
$\mathcal A\to\operatorname{Sh}(N_a)$. Their finite-limit and
colimit preservation can be checked on the conservative open
restriction family; hence the glued functor is again the inverse image
of a geometric morphism. Natural transformations glue as well. This
proves the bicategorical universal property for full topoi, not only
for local systems.

The two-open case makes passage to pro-objects completely explicit.
For a fixed target $\mathcal A$, filtered two-colimits of mapping
\emph{categories} commute with finite pseudolimits in
$\operatorname{Cat}$. This follows by representing the finitely many
objects, arrows, invertible comparison arrows, and their coherence
equalities at one common stage; the same argument handles arbitrary
morphisms and equality of morphisms. Apply this to
\[
 \operatorname{Geom}(\operatorname{Sh}(Y),\mathcal A)
 \simeq
 \operatorname{Geom}(\operatorname{Sh}(Y\setminus Z),\mathcal A)
 \mathbin{\times}^{\rm ps}_{
   \operatorname{Geom}(\operatorname{Sh}(N\setminus Z),\mathcal A)}
 \operatorname{Geom}(\operatorname{Sh}(N),\mathcal A).
\]
The left side and first factor are constant in $N$. Taking the filtered
two-colimit gives precisely the mapping-out universal property of
(F.h.7). For a pro-target $(\mathcal A_j)$, take the additional limit
over $j$; pseudolimits commute with the displayed finite pseudolimits.
This proves
(F.h.7) for arbitrary $Z$ as stated, without appealing to a putative
inverse limit of punctured full sheaf topoi. The latter can be
degenerate, as F.12 proves.

We next explain the geometric content of the extra hypothesis. On a
conical chart, compactness of $L_S$ implies that every neighbourhood of
the cone vertex contains a full radial ball of some positive radius.
Over each relatively small open subset of $S$, the same holds with a
positive local bound. A locally finite partition of unity produces a
positive continuous minorant of these bounds. Consequently tubes with
positive continuous radii form a cofinal system of all tube germs.
The control identities make these choices compatible on incident
plaques. Choose the radii successively in the finite depth, reducing
them to satisfy all earlier conditions. Only finitely many conditions
occur near a given point. Incomparable tubes are disjoint, so a
nonempty intersection of distinct tubes is indexed by an incident
flag. It is exactly a stage of (F.h.5). Recursing inside $L_S$ shows
that its angular part is the corresponding iterated-link bundle;
the remaining coordinates are positive radial intervals. This is
proved inductively on depth from (F.h.6): an incident higher plaque
is $S$ times a stratum in $L_S$ times the radial interval, and the
next tube condition is its corresponding link-tube condition. There
is no assertion that the angular bundle is a product of its fibres.

At every simultaneous radius stage, the tubes covering the given
closed union give an actual open cover of their union, and their full
intersections give its Cech nerve. Thus the preceding full-topos
descent applies. Passing from ordered intersections to all incident
flags is subdivision of the same descent datum: its one-step maps
are restrictions and its triangles are the cocycle relations. It
does not discard higher flags. For an infinite locally finite atlas,
one must also justify synchronization of pro-stages. A datum and its
isomorphisms are represented on individual charts, double overlaps,
and triple overlaps. Choose a common refinement on each such finite
family; at any point only finitely many bounds are imposed. A positive
continuous minorant gives one global radius assignment satisfying
all of them. The same argument applies to transformations and equality
of transformations. This proves essential surjectivity and full
faithfulness on the mapping categories, and justifies the pro-colimit
without commuting a filtered colimit with an arbitrary infinite
product. The same argument applies to the Cech groupoid of the stated
locally finite Hausdorff chart cover when plaques have branch
identifications: on compact chart pieces only finitely many overlap
conditions occur. This is a chart presentation of the given space,
not a claim about an arbitrary nonproper etale quotient atlas.
These constructions establish the
extension of both versions of F.13 to the stated controlled local
stratifications.

For shapes, apply the open-cover homotopy-colimit statement at every
stage. It can be proved directly by the thick Cech realization: a
partition of unity subordinate to a numerable cover provides a section
up to homotopy, and the simplex of possible partition weights over
each point is contractible. Refining charts and radii preserves this
comparison. A conical radial interval is contractible and its shrinkings
are homotopy equivalences, so the iterated-link bundles give constant
models for the resulting local \emph{pro-homotopy types}. They do not
give equivalences of full sheaf topoi: even
$\operatorname{Sh}((0,\epsilon))$ is not $\operatorname{Set}$.
The homotopy long exact sequence of each locally trivial link bundle
shows that an aspherical fibre over an aspherical base has aspherical
total space, proving the sufficient $1$-type criterion. Conversely a
space with nonzero homotopy above degree one cannot be replaced by
the classifying space of its fundamental groupoid. This is precisely
the information discarded by a level-one replacement.

The categorical cases are STRATA C.3, C.5, and C.6. The relevant
calculation in C.6 is independently checkable: for a full sieve
$A\subset C$ and complementary cosieve $B$, restriction identifies a
presheaf with $(F_A,F_B,F_B\to i^*j_*F_A)$, where
$(i^*j_*F_A)(b)=\lim_{(A/b)^{\rm op}}F_A$. The inverse sends this
triple to $j_*F_A\times_{i_*i^*j_*F_A}i_*F_B$, as is verified on
$A$ and $B$ separately. With spaces one must take the coherent limit
and pullback. Repeating this argument along a finite filtration gives
the actual crossing-flag data of C.5; its shape is the homotopy
pushout of C.3. This verifies the categorical pattern but is not
substituted for the full-topos geometric descent just proved.

Finally consider the two requested tests. For $(S^2,\{p\})$, take
$N$ a small open disc and $U=S^2\setminus\{p\}\cong\mathbb R^2$.
The punctured tube is $S^1\times(0,\epsilon)$, so (F.h.7) has shape
\[
 \operatorname{hocolim}(*\leftarrow S^1\longrightarrow *)\simeq S^2.
\]
Both open and closed strata are simply connected. The topos of sheaves
constructible for these two strata is the arrow topos: the local-system
gluing functor sends a set to that same set, since the link is connected
and the local system on $U$ is constant. Its classifying shape is
contractible. Similarly the ordinary groupoid pushout of the displayed
diagram is trivial. The full diagram with its link retains the missing
$H_2(S^2;\mathbb Z)=\mathbb Z$. Thus even aspherical elementary
pieces do not license an ordinary level-one colimit.

The drawing at scan~190 is three half-lines in the \emph{plane}:
there are seven strata, namely the origin, three open rays, and three
open sectors. The origin's link is a circle divided into three points
and three open arcs. The incident poset has seven elements, twelve
strict comparable pairs, and six strict three-element chains. Removing
the origin gives the six-vertex, six-edge incidence cycle, with first
Betti number one. Adjoining the origin cones off that cycle, as the
disc tube does. The three ray links are two-point spaces and the
sector links are empty. All these elementary angular pieces are
$1$-types, but the seven-element incident poset is not Boolean, since
seven is not a power of two. Treating the picture as only a tripod
would remove the three sectors and give a three-point link, contradicting
both the written count and the topology. The program checks the flag
counts, the Segal bijections through degree six, and the circle and
sphere homology calculations; the proof of (F.h.2) establishes all
degrees.
\end{proof}
